% preamble.tex — LaTeX header includes for the Pandoc PDF build.
% This file is injected via include-in-header, i.e. AFTER Pandoc's own
% highlighting macros, so the redefinitions of Shaded/shadecolor below
% deliberately override the defaults.

% --- Page geometry. -------------------------------------------------------
\usepackage[a4paper,top=2cm,bottom=2cm,left=2cm,right=2cm,
            includehead,headsep=8pt]{geometry}
\usepackage{microtype}

% --- Numbering: the Markdown headings already carry manual numbers (2.1,
% 5.3, ...) for the web build, so suppress LaTeX's automatic section and
% chapter numbering to avoid the doubled "2.1 2.1 ..." effect. Figures and
% tables carry their manual numbers in the caption ("Figure 2.3."), so
% LaTeX's own label is dropped as well.
\setcounter{secnumdepth}{-1}
\usepackage{caption}
\captionsetup{labelformat=empty,labelsep=none,font=small}

% --- Figures scale to the text width rather than overflow. ---------------
\usepackage{graphicx}
\setkeys{Gin}{width=0.85\linewidth,keepaspectratio}
\makeatletter
\renewcommand{\fps@figure}{htbp}
% a page of floats only starts at the top, instead of centring them with a
% blank above, and keeps 12pt between them
\setlength{\@fptop}{0pt}
\setlength{\@fpsep}{12pt}
\makeatother

% --- Breakable URLs and inline code (see break-code.lua). -----------------
\usepackage{xurl}

% --- Code blocks: wrap long lines, framed on a light background. ---------
\usepackage{fvextra}
\fvset{breaklines=true,breakanywhere=true,fontsize=\small}

\usepackage{framed}
\usepackage{xcolor}
\definecolor{shadecolor}{RGB}{244,244,248}   % very light blue-grey
\definecolor{codeframe}{gray}{0.55}
\setlength{\FrameSep}{5pt}
% (Pandoc defines Shaded only when the document has a highlighted block)
\makeatletter\@ifundefined{Shaded}{\newenvironment{Shaded}{}{}}{}\makeatother
\renewenvironment{Shaded}{%
  \def\FrameCommand{\fboxsep=\FrameSep \fcolorbox{codeframe}{shadecolor}}%
  \MakeFramed{\advance\hsize-\width \FrameRestore}}%
 {\endMakeFramed}

% --- A discreet running header. ------------------------------------------
\usepackage{fancyhdr}
\pagestyle{fancy}
\fancyhf{}
\fancyhead[L]{\small PBeFS — Parallel Best-First Search}
\fancyhead[R]{\small\thepage}
\renewcommand{\headrulewidth}{0.2pt}
\setlength{\headheight}{14pt}

% --- The Unicode symbols of the pseudocode (the same as in the Doxygen
% comments of the headers) and of the prose, mapped so that pdflatex
% typesets them in text and verbatim alike.
\usepackage{amssymb}
\DeclareUnicodeCharacter{2264}{\ensuremath{\le}}         % <=
\DeclareUnicodeCharacter{2265}{\ensuremath{\ge}}         % >=
\DeclareUnicodeCharacter{2260}{\ensuremath{\ne}}         % not equal
\DeclareUnicodeCharacter{221E}{\ensuremath{\infty}}      % infinity
\DeclareUnicodeCharacter{2208}{\ensuremath{\in}}         % element of
\DeclareUnicodeCharacter{2209}{\ensuremath{\notin}}      % not element of
\DeclareUnicodeCharacter{220B}{\ensuremath{\ni}}         % contains
\DeclareUnicodeCharacter{2205}{\ensuremath{\emptyset}}   % empty set
\DeclareUnicodeCharacter{222A}{\ensuremath{\cup}}        % union
\DeclareUnicodeCharacter{2229}{\ensuremath{\cap}}        % intersection
\DeclareUnicodeCharacter{2216}{\ensuremath{\setminus}}   % set minus
\DeclareUnicodeCharacter{2286}{\ensuremath{\subseteq}}   % subset
\DeclareUnicodeCharacter{2190}{\ensuremath{\leftarrow}}  % <-
\DeclareUnicodeCharacter{2192}{\ensuremath{\rightarrow}} % ->
\DeclareUnicodeCharacter{25B7}{\ensuremath{\triangleright}} % comment
\DeclareUnicodeCharacter{2212}{\ensuremath{-}}           % minus sign
\DeclareUnicodeCharacter{00B7}{\ensuremath{\cdot}}       % middle dot
\DeclareUnicodeCharacter{00D7}{\ensuremath{\times}}      % times
\DeclareUnicodeCharacter{2026}{\ldots}                   % ellipsis
\DeclareUnicodeCharacter{25CF}{\textbullet}              % black circle
\DeclareUnicodeCharacter{220E}{\ensuremath{\square}}     % end of proof
\DeclareUnicodeCharacter{2080}{\ensuremath{_0}}          % subscript 0
\DeclareUnicodeCharacter{2081}{\ensuremath{_1}}          % subscript 1
\DeclareUnicodeCharacter{2082}{\ensuremath{_2}}          % subscript 2
\DeclareUnicodeCharacter{1D62}{\ensuremath{_i}}          % subscript i
\DeclareUnicodeCharacter{2C7C}{\ensuremath{_j}}          % subscript j
\DeclareUnicodeCharacter{03B1}{\ensuremath{\alpha}}
\DeclareUnicodeCharacter{03B2}{\ensuremath{\beta}}
\DeclareUnicodeCharacter{03B3}{\ensuremath{\gamma}}
\DeclareUnicodeCharacter{03B4}{\ensuremath{\delta}}
\DeclareUnicodeCharacter{03B5}{\ensuremath{\varepsilon}}
\DeclareUnicodeCharacter{03BB}{\ensuremath{\lambda}}
\DeclareUnicodeCharacter{03C0}{\ensuremath{\pi}}
\DeclareUnicodeCharacter{03C3}{\ensuremath{\sigma}}
\DeclareUnicodeCharacter{03C4}{\ensuremath{\tau}}
\DeclareUnicodeCharacter{03A3}{\ensuremath{\Sigma}}
